%% toc: P2 21 21.8
\section{B\"uchi's prefix form of the refined reduction formula}
\label{sec:p02-w02c-p02-025-buchi-form}

\emph{Purpose.}  The refined reduction formula \(\rho_{M,\prec}\) of
\cref{sec:p02-w02-p02-036-refined} carries Trakhtenbrot's theorem in its
inseparable form, but it spends a binary successor symbol \(S\), one binary
tape symbol for every tape letter, and the prefix
\(\forall x\,\exists u\,\exists v\,\forall y\) in its relativized infinity
axiom.  The source's Exercise~2.1.31 asks for much less
\cite[Sect.~2.1.2, Exercise~2.1.31, pp.~35--36]{BorgerGradelGurevich1997}:
show that \(\rho_{M,\prec}\) can be chosen as a formula of the form
\[
 \exists x\,Zx\;\land\;\forall x\,\exists u\,\forall y\,\beta ,
\]
where \(Z\) is a monadic predicate symbol and \(\beta\) is a formula whose only
variables are \(x\), \(u\) and \(y\) and whose only predicate symbols are the
three binary symbols \(H\), \(T_0\), \(K\) and finitely many monadic ones.  The
source's hint is to use Turing machines with the binary alphabet \(\{0,1\}\),
to read \(Zx\) as \(x=0\), and to place \(\mathrm{START}\) inside \(\beta\) as
\(Zy\to\mathrm{START}[0/y]\).  This section carries the exercise out in full,
with \(\beta\) quantifier-free and without equality, and proves the refined
reduction property for the resulting formulas.  It also proves the remark that
follows the exercise in the source: every model of such a formula satisfies
\(\exists y\,Hyy\).

\emph{Inputs.}  The Turing machines of \cref{sec:turing-machines-01}, with
their two-way tape, and their computations on the empty input;
the program codes \(e\), the decidable set \(V\) of valid codes and the decoded
machines \(M_e\) of \cref{sec:universal-reductions-01}; the halting sets
\(H_1\), \(H_2\), \(H_\infty\) of
Definition~\ref{def:p02-w02-p02-036-halting}; the classes
\(\mathrm{Fin\text{-}sat}\), \(\mathrm{Non\text{-}sat}\) and
\(\mathrm{Inf\text{-}axioms}\) of
Definition~\ref{def:p02-w02-p02-036-classes}; and recursive inseparability
(Definition~\ref{def:p02-w02-p02-036-inseparable}).  Structures have nonempty
domains.  No formula below contains the equality sign.

\emph{Conclusions.}  For every valid code \(e\) there is a sentence
\(\rho^{\mathrm B}_{e}\) of the displayed form, computable from \(e\), such that
\(M_e\) on the empty input rejects exactly when \(\rho^{\mathrm B}_e\) is
unsatisfiable, accepts exactly when \(\rho^{\mathrm B}_e\) has a finite model, and
diverges exactly when \(\rho^{\mathrm B}_e\) is an infinity axiom
(\cref{thm:p02-w02c-p02-025-exercise}).  Consequently the finitely
satisfiable, the unsatisfiable and the only infinitely satisfiable sentences of
this prefix form are pairwise recursively inseparable
(\cref{thm:p02-w02c-p02-025-inseparable}).  The route passes through binary
one-way programs that move slowly
(Definition~\ref{def:p02-w02c-p02-025-program}), to which every Turing
machine reduces with its halting behaviour preserved
(\cref{thm:p02-w02c-p02-025-universality}).  All proofs are in the annex.

\subsection{Binary one-way programs}
\label{sec:p02-w02c-p02-025-programs}

The source's machines in Sect.~2.1 have a one-way infinite tape with cells
\(0,1,2,\ldots\), may leave the head where it is, and start in state \(0\) at
the left end of the empty tape; for the halting problems \(H_1\) and \(H_2\)
they have exactly the two stop states \(1\) and \(2\)
\cite[Sect.~2.1.1, p.~20; Sect.~2.1.2, p.~34]{BorgerGradelGurevich1997}.  The
next definition fixes that format for the binary alphabet and adds one timing
condition that the formula below needs.

\begin{hbdefinition}{def:p02-w02c-p02-025-program}{Binary one-way program and admissibility}
A \emph{binary one-way program} is a pair \(P=(r,\delta_P)\) with \(r\geq2\) and
\[
 \delta_P:\bigl(\{0,\ldots,r\}\setminus\{1,2\}\bigr)\times\{0,1\}
 \longrightarrow\{0,\ldots,r\}\times\{0,1\}\times\{L,N,R\}.
\]
Its states are \(0,\ldots,r\); state \(0\) is the start state and \(1\), \(2\)
are the stop states.  The letter \(0\) is the blank.  A \emph{configuration} is
a triple \((p,B,h)\) with \(p\) a state, \(h\in\mathbb N\) and
\(B:\mathbb N\to\{0,1\}\) zero outside a finite set.  If \(p\notin\{1,2\}\) and
\(\delta_P(p,B(h))=(p',b,d)\), put \(B'\defeq B[h\mapsto b]\) and
\(h'\defeq h-1,h,h+1\) for \(d=L,N,R\); when \(h'\geq0\) the configuration
\((p,B,h)\) has the successor \((p',B',h')\).  A configuration in a stop state,
and one whose step would move the head left from cell \(0\), has no
successor.  The \emph{computation} of \(P\) is the maximal sequence
\(c_0,c_1,\ldots\) of configurations with \(c_0=(0,\bar0,0)\), where
\(\bar0\) is the all-zero tape, and \(c_{t+1}\) the successor of \(c_t\).  Write
\(c_t=(q_t,B_t,h_t)\) and let \(d_t\) be the move of step \(t\).  \(P\)
\emph{halts in state \(i\)} (\(i=1,2\)) when some \(q_s=i\).

\(P\) is \emph{admissible} when its computation has the following two
properties:
\begin{enumerate}
\item[(N1)] no configuration of the computation fails to have a successor
  except one in a stop state, that is, the head never moves left from cell
  \(0\);
\item[(N2)] whenever \(d_t=R\), then \(h_t\leq t-2\).
\end{enumerate}
\index{binary one-way program}\index{admissible program}
\end{hbdefinition}

The letters of the tape are \(0\) and \(1\) only; the symbol \(T_0\) of the
formula below will record the letter \(0\), and its negation the letter \(1\).
Condition (N2) says that the head travels to the right only well behind the
current time; it is what lets one binary symbol hold two kinds of information
on the two sides of the diagonal.

\begin{hblemma}{thm:p02-w02c-p02-025-admissible}{proved=annex:p02-w02c-p02-025-admissible}{Elementary bounds for admissible programs}
Let \(P\) be admissible with computation \((q_t,B_t,h_t)_t\).  For every \(t\)
for which \(c_t\) exists: \(h_t\leq t\); \(h_t<t\) if \(t\geq1\); and
\(B_t(c)=0\) for every cell \(c\geq t\).
\end{hblemma}

\begin{hbdefinition}{def:p02-w02c-p02-025-simulation}{Block simulation of a Turing machine}
Let \(M=(Q,\mathcal A_{\mathrm{in}},\mathcal A,\delta,q_0,q_{\mathrm{acc}},
q_{\mathrm{rej}})\) be a Turing machine with tape alphabet \(\mathcal A\), and
put \(Q^{\circ}\defeq Q\setminus\{q_{\mathrm{acc}},q_{\mathrm{rej}}\}\).  Choose
the least \(m\geq1\) with \(2^m\geq\lvert\mathcal A\rvert\), an injective
\(\gamma:\mathcal A\to\{0,1\}^m\) with \(\gamma(\sqcup)=0^m\), and put
\(k\defeq2m+1\).  For a two-way tape \(T:\mathbb Z\to\mathcal A\) and \(j\geq0\),
the \emph{block} of \(T\) at \(j\) is the word
\(\operatorname{blk}_T(j)\defeq f_j\,\gamma(T(j))\,\gamma(T(-j-1))\in\{0,1\}^k\),
where \(f_j=1\) if \(j=0\) and \(f_j=0\) otherwise; the binary tape
\(\operatorname{enc}(T)\) has \(\operatorname{enc}(T)(kj+i)\) equal to the
letter number \(i\) (counting from \(0\)) of \(\operatorname{blk}_T(j)\).
For \(h\in\mathbb Z\) put \(\operatorname{pos}(h)\defeq(h,+)\) if \(h\geq0\)
and \(\operatorname{pos}(h)\defeq(-h-1,-)\) if \(h<0\).  For a word
\(v\in\{0,1\}^k\) write \(v_i\) for its letter number \(i\), and call \(v=f\,x\,y\)
with \(\lvert x\rvert=\lvert y\rvert=m\) \emph{proper} if
\(x,y\in\gamma[\mathcal A]\).

The program \(\operatorname{Sim}(M)\) has the states
\(\iota\) (start), \(\bot_1\), \(\bot_2\) (stop), \(\lambda\), and
\begin{itemize}
\item \(\mathrm{rd}(q,\sigma,w)\) for \(q\in Q^{\circ}\), \(\sigma\in\{+,-\}\)
  and \(w\in\{0,1\}^{<k}\);
\item \(\mathrm{wr}(v,i,q',\sigma',\Delta)\) for \(v\in\{0,1\}^k\),
  \(0\leq i\leq k-2\), \(q'\in Q^{\circ}\), \(\sigma'\in\{+,-\}\) and
  \(\Delta\in\{-1,0,1\}\);
\item \(\mathrm{tr}(D,q',\sigma',n)\) for \(D\in\{L,R\}\),
  \(1\leq n\leq k-1\), \(q'\in Q^{\circ}\) and \(\sigma'\in\{+,-\}\).
\end{itemize}
Its transitions, for each letter \(a\in\{0,1\}\), are:
\begin{enumerate}
\item \(\iota\): to \((\bot_1,1,N)\) if \(q_0=q_{\mathrm{rej}}\), to
  \((\bot_2,1,N)\) if \(q_0=q_{\mathrm{acc}}\), and otherwise to
  \((\mathrm{rd}(q_0,+,\varepsilon),1,N)\).
\item \(\mathrm{rd}(q,\sigma,w)\) with \(\lvert w\rvert<k-1\): to
  \((\mathrm{rd}(q,\sigma,wa),a,R)\).
\item \(\mathrm{rd}(q,\sigma,w)\) with \(\lvert w\rvert=k-1\): let
  \(v\defeq wa=f\,x\,y\).  If \(v\) is not proper, to \((\lambda,a,N)\).
  Otherwise let \(a^{*}\defeq\gamma^{-1}(x)\) if \(\sigma=+\) and
  \(a^{*}\defeq\gamma^{-1}(y)\) if \(\sigma=-\), and let
  \(\delta(q,a^{*})=(q',b,d)\).  If \(q'=q_{\mathrm{rej}}\), to
  \((\bot_1,a,N)\); if \(q'=q_{\mathrm{acc}}\), to \((\bot_2,a,N)\).  Otherwise
  let \(v'\defeq f\,\gamma(b)\,y\) if \(\sigma=+\) and
  \(v'\defeq f\,x\,\gamma(b)\) if \(\sigma=-\), let \((\Delta,\sigma')\) be
  given by the table
  \[
  \begin{array}{c|c|c|c}
   \sigma & d & f & (\Delta,\sigma')\\\hline
   + & R & 0,1 & (+1,+)\\
   + & L & 0 & (-1,+)\\
   + & L & 1 & (0,-)
  \end{array}
  \qquad
  \begin{array}{c|c|c|c}
   \sigma & d & f & (\Delta,\sigma')\\\hline
   - & L & 0,1 & (+1,-)\\
   - & R & 0 & (-1,-)\\
   - & R & 1 & (0,+)
  \end{array}
  \]
  and go to \((\mathrm{wr}(v',k-2,q',\sigma',\Delta),v'_{k-1},L)\).
\item \(\mathrm{wr}(v',i,q',\sigma',\Delta)\) with \(i\geq1\): to
  \((\mathrm{wr}(v',i-1,q',\sigma',\Delta),v'_i,L)\).
\item \(\mathrm{wr}(v',0,q',\sigma',\Delta)\): to
  \((\mathrm{rd}(q',\sigma',\varepsilon),v'_0,N)\) if \(\Delta=0\), to
  \((\mathrm{tr}(R,q',\sigma',k-1),v'_0,R)\) if \(\Delta=+1\), and to
  \((\mathrm{tr}(L,q',\sigma',k-1),v'_0,L)\) if \(\Delta=-1\).
\item \(\mathrm{tr}(D,q',\sigma',n)\) with \(n\geq2\): to
  \((\mathrm{tr}(D,q',\sigma',n-1),a,D)\).
\item \(\mathrm{tr}(D,q',\sigma',1)\): to
  \((\mathrm{rd}(q',\sigma',\varepsilon),a,D)\).
\item \(\lambda\): to \((\lambda,a,N)\).
\end{enumerate}
The program \(\operatorname{Slow}(P')\) of a program \(P'\) with start state
\(\iota'\) and stop states \(\bot_1,\bot_2\) has the states \(a_0\) (start),
\(a_1\), \(\bot_1\), \(\bot_2\), and \((p,0)\), \((p,1)\) for every non-stop
state \(p\) of \(P'\), with the transitions, for \(a\in\{0,1\}\):
\(a_0\mapsto(a_1,a,N)\); \(a_1\mapsto((\iota',0),a,N)\);
\((p,1)\mapsto((p,0),a,N)\); and, if \(\delta_{P'}(p,a)=(p',b,d)\),
\((p,0)\mapsto(p',b,d)\) when \(p'\) is a stop state and
\((p,0)\mapsto((p',1),b,d)\) otherwise.  For a valid code \(e\), the program
\(P_e\) is \(\operatorname{Slow}(\operatorname{Sim}(M_e))\) with its states
numbered \(0,\ldots,r\) in a fixed computable order that sends \(a_0\) to
\(0\), \(\bot_1\) to \(1\) and \(\bot_2\) to \(2\).  Configurations, successor
steps and computations of \(\operatorname{Sim}(M)\) and of
\(\operatorname{Slow}(P')\) are defined as in
Definition~\ref{def:p02-w02c-p02-025-program}, with the named start and stop
states in place of \(0\), \(1\), \(2\).
\index{block simulation}
\end{hbdefinition}

The program \(\operatorname{Sim}(M)\) folds the two-way tape: block \(j\) holds
the cells \(j\) and \(-j-1\) of \(M\) in two tracks, the flag \(f_0=1\) marks
the block where the fold turns, and the sign \(\sigma\) in the state records
the track on which \(M\)'s head stands.  One step of \(M\) costs \(k\) reading
steps, \(k-1\) writing steps and, unless the head stays in its block,
\(k-1\) further travelling steps.  The program \(\operatorname{Slow}(P')\) inserts two idle
steps at the start and one idle step after each step of \(P'\), so that every
move happens at an even time \(t\geq2\).

\begin{hbtheorem}{thm:p02-w02c-p02-025-universality}{proved=annex:p02-w02c-p02-025-universality}{Binary admissible programs simulate Turing machines}
For every valid code \(e\), the program \(P_e\) is admissible and is computable
from \(e\), and on the empty input:
\begin{enumerate}
\item \(M_e\) rejects if and only if \(P_e\) halts in state \(1\);
\item \(M_e\) accepts if and only if \(P_e\) halts in state \(2\);
\item \(M_e\) diverges if and only if the computation of \(P_e\) is infinite.
\end{enumerate}
Hence \(e\in H_1\), \(e\in H_2\), \(e\in H_\infty\) according as \(P_e\) halts
in state \(1\), halts in state \(2\), or never halts.
\end{hbtheorem}

This is the universality of binary-alphabet machines that the source's hint
invokes, in the form needed here: it keeps apart the two stop states and the
non-halting case, it uses the one-way tape and the empty start, and it adds the
timing condition (N2).

\subsection{The formula}
\label{sec:p02-w02c-p02-025-formula}

\begin{hbdefinition}{def:p02-w02c-p02-025-formula}{The formula \(\rho^{\mathrm B}_{P}\)}
Let \(P=(r,\delta_P)\) be a binary one-way program.  The signature
\(\Lang^{\mathrm B}_P\) consists of the binary relation symbols \(H\), \(T_0\),
\(K\) and the monadic relation symbols \(Z\), \(I_0,\ldots,I_r\).  Write
\([0]vw\defeq T_0vw\) and \([1]vw\defeq\neg T_0vw\) for variables \(v,w\), and
\(\mathrm{Act}\,x\defeq\bigvee_{p\notin\{1,2\}}I_px\).  Let \(\beta_P\) be the
conjunction of the following quantifier-free formulas in the variables
\(x,u,y\).
\begin{enumerate}
\item[(\(\beta\)1)] \(\mathrm{START}\):\quad
  \(Zy\to\bigl(I_0y\land Hyy\land T_0yx\bigr)\).
\item[(\(\beta\)2)] \(\mathrm{BLANK}\):\quad \(T_0xx\).
\item[(\(\beta\)3)] \(\mathrm{OFFDIAG}\):\quad \((Zx\land Kxy)\to\neg Hyy\).
\item[(\(\beta\)4)] \(\mathrm{STEP}_P\): for each state \(p\notin\{1,2\}\) and
  letter \(a\) with \(\delta_P(p,a)=(p',b,d)\),
  \[
   \bigl(Hxy\land\neg Kxy\land I_px\land[a]xy\bigr)\to
   \bigl(I_{p'}u\land[b]uy\bigr),
  \]
  and, according to \(d\),
  \[
  \begin{aligned}
   d=N:&\quad\bigl(Hxy\land\neg Kxy\land I_px\land[a]xy\bigr)\to Huy,\\
   d=R:&\quad\bigl(Hyx\land\neg Kyx\land I_py\land[a]yx\bigr)\to Huy,\\
   d=L:&\quad\bigl(Hyu\land\neg Kyu\land I_py\land[a]yu\bigr)\to Hxy.
  \end{aligned}
  \]
\item[(\(\beta\)5)] \(\mathrm{MOVE}\):\quad \((Hyx\land Kyx)\to Huy\).
\item[(\(\beta\)6)] \(\mathrm{FRAME}\):\quad
  \((\neg Kxy\land\neg Hxy\land T_0xy)\to T_0uy\) and
  \((\neg Kxy\land\neg Hxy\land\neg T_0xy)\to\neg T_0uy\).
\item[(\(\beta\)7)] \(\mathrm{ONEHEAD}\): the three formulas
  \((Hyx\land\neg Kyx)\to T_0xy\), \((T_0uy\land Kuy)\to T_0xy\) and
  \(\neg(T_0uy\land Kuy\land Hyx)\).
\item[(\(\beta\)8)] \(\mathrm{NONSTOP}\):\quad \(\neg I_1x\).
\item[(\(\beta\)9)] \(\prec_{K,P}\) and activity:
  \[
   \neg Kxx\land(Hyx\to Kxu)\land(Kyx\to Kyu)\land
   (\mathrm{Act}\,x\to Kxu).
  \]
\end{enumerate}
Put
\[
 \rho^{\mathrm B}_P\;\defeq\;\exists x\,Zx\;\land\;\forall x\,\exists u\,
 \forall y\,\beta_P ,
\]
and, for a valid code \(e\), \(\rho^{\mathrm B}_e\defeq\rho^{\mathrm B}_{P_e}\)
with \(P_e\) as in Definition~\ref{def:p02-w02c-p02-025-simulation}.
\index{reduction formula!B\"uchi form}
\end{hbdefinition}

The intended model has as domain an initial segment of \(\mathbb N\), reads
\(Zx\) as \(x=0\), \(Kxy\) as \(x<y\), and takes \(u\) to be \(x+1\).  One
element serves both as a moment of time and as a tape cell, and the
order \(K\) splits the pairs \((v,w)\) into three regions, each carrying its own
information.
\begin{itemize}
\item Below the diagonal (\(w<v\)): \(Hvw\) says that at time \(v\) the head is
  on cell \(w\); \(T_0vw\) says that at time \(v\) cell \(w\) holds \(0\).
\item On the diagonal (\(w=v\)): \(Hvv\) says that at time \(v\) the head is on
  cell \(v\), which happens only for \(v=0\); \(T_0vv\) always holds, since cell
  \(v\) is still blank at time \(v\).
\item Above the diagonal (\(v<w\)): \(Hvw\) says that the step at time \(w\)
  moves the head and leaves it on cell \(v\) at time \(w+1\); \(T_0vw\) says
  that \(v\) is at most the head position at time \(w\).
\end{itemize}
By Lemma~\ref{thm:p02-w02c-p02-025-admissible}, a cell at or to the right of
the current time is still blank and the head stands strictly left of the
current time after time \(0\); so the lower region and the diagonal hold the
whole configuration, and the upper region is free for two auxiliary records.
The conjuncts read as follows.  \(\mathrm{START}\) is the source's
\(\mathrm{START}=I_0(0)\land H(0,0)\land\forall x\,T_0(0,x)\)
\cite[Sect.~2.1.1, p.~22]{BorgerGradelGurevich1997} with \(0\) replaced by
\(y\) under the hypothesis \(Zy\); its quantifier \(\forall x\) is supplied
by the outer \(\forall x\) of \(\rho^{\mathrm B}_P\), since the conjunct does
not mention \(u\).  In the upper region, \(T_0(0,w)\) says that cell \(0\) is
not to the right of the head at time \(w\), which is true.  \(\mathrm{STEP}_P\)
takes the place of the source's step formulas (2.1)--(2.3)
\cite[Sect.~2.1.1, pp.~20--21]{BorgerGradelGurevich1997}, with the successor
relation \(S\) replaced by the witness \(u\) and the letter \(1\) written as
\(\neg T_0\).  A move is made in two stages because a
single instance of \(\beta_P\) sees the pairs \((x,y)\), \((u,y)\), \((y,x)\)
and \((y,u)\) but never a pair whose both entries are successors: the head
position after a move is first written into the upper region by
\(\mathrm{STEP}_P\) (cell \(x+1\) or cell \(u-1\), at time \(y\)) and then
copied down into the next row by \(\mathrm{MOVE}\).  \(\mathrm{FRAME}\) is
the source's ``no change outside the head cell'' (2.3) with the unavailable
inequality \(y\neq x\) replaced by \(\neg Hxy\).  \(\mathrm{ONEHEAD}\) marks, at
each time, the cells at or left of a head and forbids a head with a mark to its
right; together with \(\mathrm{OFFDIAG}\) it makes the head unique, which is
what makes \(\neg Hxy\) in \(\mathrm{FRAME}\) mean ``not the head cell''.
\(\mathrm{NONSTOP}\) is the source's.  The first three conjuncts of
(\(\beta\)9) are the source's relativized infinity axiom
\(\prec_{K,M}=\forall x\,\exists u\,\forall y\,(\neg Kxx\land(Hyx\to Kxu)\land
(Kyx\to Kyu))\) \cite[Sect.~2.1.2, p.~35]{BorgerGradelGurevich1997}; the
added \(\mathrm{Act}\,x\to Kxu\) extends the existential claim from visited
cells to the moments before the program halts, because here the moments are
elements of the same order.

\begin{hblemma}{thm:p02-w02c-p02-025-shape}{proved=annex:p02-w02c-p02-025-shape}{Shape of \(\rho^{\mathrm B}_P\)}
For every binary one-way program \(P\), the sentence \(\rho^{\mathrm B}_P\) has
the form \(\exists x\,Zx\land\forall x\,\exists u\,\forall y\,\beta_P\) with
\(Z\) monadic and \(\beta_P\) quantifier-free; the only variables of
\(\beta_P\) are \(x\), \(u\), \(y\); the only binary predicate symbols in
\(\rho^{\mathrm B}_P\) are \(H\), \(T_0\), \(K\), and the others are the
\(r+2\) monadic symbols \(Z,I_0,\ldots,I_r\); equality does not occur.  The
map \(P\mapsto\rho^{\mathrm B}_P\) is computable, and so is
\(e\mapsto\rho^{\mathrm B}_e\) on \(V\).
\end{hblemma}

\begin{hblemma}{thm:p02-w02c-p02-025-remark}{proved=annex:p02-w02c-p02-025-remark}{The remark: every model satisfies \(\exists y\,Hyy\)}
For every binary one-way program \(P\), every model of \(\rho^{\mathrm B}_P\)
satisfies \(\exists y\,Hyy\).
\end{hblemma}

The source states this remark after the exercise, for use in its next chapter
\cite[Sect.~2.1.2, p.~36]{BorgerGradelGurevich1997}, and attributes it to the
subformula that makes the initial head position diagonal; in \(\beta_P\) that
subformula is the conjunct \(Zy\to Hyy\) of \(\mathrm{START}\).  The next
chapter of the source uses it in Step~1 of its proof of the
Kalm\'ar--Sur\'anyi theorem, where formulas
\(\exists v\,\alpha\land\forall x\,\exists u\,\forall y\,\beta\) of this kind
are called B\"uchi formulas and the remark is quoted as giving \(\exists v\,Hvv\)
in every model \cite[Sect.~3.2, pp.~115--116]{BorgerGradelGurevich1997}.

\subsection{Simulation and the refined reduction property}
\label{sec:p02-w02c-p02-025-reduction}

\begin{hbdefinition}{def:p02-w02c-p02-025-intended}{The intended structure}
Let \(P\) be an admissible program whose computation
\((q_t,B_t,h_t)_t\) does not reach state \(1\).  Let \(D\defeq\{0,\ldots,s\}\)
if it halts in state \(2\) at time \(s\), and \(D\defeq\mathbb N\) if it never
halts.  The structure \(\mathfrak A_P\) has domain \(D\) and
\[
\begin{aligned}
 Z&=\{0\},\qquad K=\{(v,w)\in D^2:v<w\},\qquad I_p=\{t\in D:q_t=p\},\\
 H&=\{(t,h_t):t\in D\}\cup\{(h_{t+1},t):t+1\in D,\ d_t\neq N\},\\
 T_0&=\{(t,c):c<t,\ B_t(c)=0\}\cup\{(t,t):t\in D\}\cup
     \{(c,t):c<t,\ c\leq h_t\},
\end{aligned}
\]
all pairs ranging over \(D\).  Its witness map is \(g(v)\defeq v+1\) if
\(v+1\in D\) and \(g(v)\defeq v\) otherwise.
\end{hbdefinition}

\begin{hblemma}{thm:p02-w02c-p02-025-intended}{proved=annex:p02-w02c-p02-025-intended}{The intended structure is a model}
In the setting of Definition~\ref{def:p02-w02c-p02-025-intended},
\(\mathfrak A_P\entails\beta_P[v,g(v),w]\) for all \(v,w\in D\), so
\(\mathfrak A_P\entails\rho^{\mathrm B}_P\).  If the computation halts at time
\(s\), the model has exactly \(s+1\) elements.
\end{hblemma}

\begin{hblemma}{thm:p02-w02c-p02-025-simulation}{proved=annex:p02-w02c-p02-025-simulation}{Simulation lemma for \(\rho^{\mathrm B}_P\)}
Let \(P\) be admissible with computation \((q_t,B_t,h_t)_{t<\ell}\), where
\(\ell\in\mathbb N\cup\{\infty\}\) is its length, let
\(\mathfrak A\entails\rho^{\mathrm B}_P\), and let \(z\in Z^{\mathfrak A}\).
Then there are elements \(\tau_0=z,\tau_1,\tau_2,\ldots\) of \(\mathfrak A\)
such that \(\neg K^{\mathfrak A}(\tau_b,\tau_a)\) for all \(a\leq b\), and for
every \(t<\ell\):
\begin{enumerate}
\item \(\mathfrak A\entails I_{q_t}[\tau_t]\);
\item \(K^{\mathfrak A}(\tau_a,\tau_b)\) for all \(a<b\leq t\); in particular
      \(\tau_0,\ldots,\tau_t\) are pairwise distinct;
\item \(\mathfrak A\entails H[\tau_t,\tau_{h_t}]\), and
      \(\mathfrak A\entails\neg H[\tau_t,\tau_c]\) for every \(c\leq t\) with
      \(c\neq h_t\);
\item for every \(c\leq t\), \(\mathfrak A\entails T_0[\tau_t,\tau_c]\) if and
      only if \(B_t(c)=0\).
\end{enumerate}
\end{hblemma}

The lemma is the source's Simulation Lemma, its Lemma~2.1.2
\cite[Sect.~2.1.1, p.~20]{BorgerGradelGurevich1997}, in this setting, with one difference forced by the alphabet.  The source's formulas are
frugal: a model may contain more than the intended facts, for instance two
letters in one cell.  With the letter \(1\) written as \(\neg T_0\), a cell
holds exactly one letter, so the model must agree with the computation exactly
on the tape; the uniqueness of the head in item~3 is what secures this.

\begin{hbtheorem}{thm:p02-w02c-p02-025-reduction}{proved=annex:p02-w02c-p02-025-reduction}{Refined reduction property for \(\rho^{\mathrm B}_P\)}
For every admissible binary one-way program \(P\):
\begin{enumerate}
\item \(P\) halts in state \(1\) if and only if \(\rho^{\mathrm B}_P\) is
      unsatisfiable;
\item \(P\) halts in state \(2\) if and only if \(\rho^{\mathrm B}_P\) has a
      finite model; if it halts at time \(s\), \(\rho^{\mathrm B}_P\) has a model
      with exactly \(s+1\) elements;
\item \(P\) never halts if and only if \(\rho^{\mathrm B}_P\) is satisfiable
      and has only infinite models.
\end{enumerate}
\end{hbtheorem}

\begin{hbtheorem}{thm:p02-w02c-p02-025-exercise}{proved=annex:p02-w02c-p02-025-exercise}{The source's Exercise~2.1.31}
The refined reduction formula can be chosen in the form
\(\exists x\,Zx\land\forall x\,\exists u\,\forall y\,\beta\) with \(Z\)
monadic, \(\beta\) quantifier-free in the variables \(x,u,y\), and only the
binary predicate symbols \(H\), \(T_0\), \(K\) besides finitely many monadic
ones: for every valid code \(e\), the sentence \(\rho^{\mathrm B}_e\) has this
form, \(e\mapsto\rho^{\mathrm B}_e\) is computable on \(V\), and
\begin{enumerate}
\item \(e\in H_1\) if and only if \(\rho^{\mathrm B}_e\) is contradictory;
\item \(e\in H_2\) if and only if \(\rho^{\mathrm B}_e\) is finitely
      satisfiable;
\item \(e\in H_\infty\) if and only if
      \(\rho^{\mathrm B}_e\in\mathrm{Inf\text{-}axioms}\).
\end{enumerate}
\end{hbtheorem}

Items 1 and~2 are the source's reduction property (i) and (ii) for
\(\rho_{M,\prec}\) \cite[Sect.~2.1.2, p.~34]{BorgerGradelGurevich1997}, now for
formulas of the prefix form the exercise asks for; item~3 follows from them
because \(H_1\), \(H_2\), \(H_\infty\) partition \(V\).

\begin{hbcorollary}{thm:p02-w02c-p02-025-inseparable}{proved=annex:p02-w02c-p02-025-inseparable}{Trakhtenbrot's theorem for the B\"uchi form}
Let \(\mathcal B\) be the set of sentences
\(\exists x\,Zx\land\forall x\,\exists u\,\forall y\,\beta\) with \(\beta\)
quantifier-free in the variables \(x,u,y\), built without equality from the
binary symbols \(H\), \(T_0\), \(K\) and monadic symbols from the list
\(Z,I_0,I_1,\ldots\).  Then \(\mathrm{Fin\text{-}sat}(\mathcal B)\),
\(\mathrm{Non\text{-}sat}(\mathcal B)\) and
\(\mathrm{Inf\text{-}axioms}(\mathcal B)\) are pairwise recursively
inseparable.
\end{hbcorollary}

In the source's notation for prefix-vocabulary classes this is the class
\([\exists\land\forall\exists\forall,(\omega,3)]\) of B\"uchi formulas with
monadic and three binary predicates
\cite[Sect.~3.2, p.~115]{BorgerGradelGurevich1997}.  The corollary does not
say that this class is a conservative reduction class in the sense of
Definition~\ref{def:p02-w02-p02-036-conservative}; that stronger property,
which the source asserts for it on the same page, needs a reduction of all
first-order formulas, not only of the three halting problems.

\subsection{Worked example}
\label{sec:p02-w02c-p02-025-example}

\begin{hbexample}{def:p02-w02c-p02-025-example}{A three-step program}
Let \(r=4\) and let \(P\) have
\[
\begin{aligned}
 \delta_P(0,a)&=(3,1,N), & \delta_P(3,a)&=(4,1,N),\\
 \delta_P(4,1)&=(2,0,R), & \delta_P(4,0)&=(4,0,N)
\end{aligned}
\]
for \(a\in\{0,1\}\).  Its computation is
\[
 (0,\bar0,0),\quad(3,1\bar0,0),\quad(4,1\bar0,0),\quad(2,\bar0,1),
\]
writing \(1\bar0\) for the tape with \(1\) in cell \(0\) and \(0\) elsewhere.
It halts in state \(2\) at time \(s=3\).  The only move is \(d_2=R\) with
\(h_2=0\leq 2-2\), so \(P\) is admissible.  The intended model of
Definition~\ref{def:p02-w02c-p02-025-intended} has domain \(\{0,1,2,3\}\),
witness map \(0\mapsto1\mapsto2\mapsto3\mapsto3\), \(K={<}\), \(Z=\{0\}\),
\(I_0=\{0\}\), \(I_3=\{1\}\), \(I_4=\{2\}\), \(I_2=\{3\}\), \(I_1=\varnothing\), and
\[
\begin{aligned}
 H&=\{(0,0),(1,0),(2,0),(3,1)\}\cup\{(1,2)\},\\
 T_0&=\{(2,1),(3,0),(3,1),(3,2)\}\cup\{(0,0),(1,1),(2,2),(3,3)\}\\
   &\qquad\cup\{(0,1),(0,2),(0,3),(1,3)\}.
\end{aligned}
\]
The pair \((1,2)\) of \(H\) is the upper record of the move at time \(2\): the
instance \(x=0\), \(u=1\), \(y=2\) of the \(d=R\) clause has the premises
\(H(2,0)\), \(\neg K(2,0)\), \(I_4(2)\) and \(\neg T_0(2,0)\), and yields
\(H(1,2)\); the instance \(x=2\), \(u=3\), \(y=1\) of \(\mathrm{MOVE}\) then
yields the head position \(H(3,1)\).  The pairs \((1,0)\) and \((2,0)\) are
missing from \(T_0\) because cell \(0\) holds \(1\) at times \(1\) and \(2\).
The upper pairs of \(T_0\) are the marks of \(\mathrm{ONEHEAD}\): at time \(3\)
the head is on cell \(1\), so cells \(0\) and \(1\) are marked,
\((0,3),(1,3)\in T_0\), and \((2,3)\notin T_0\).  The model has
\(s+1=4\) elements, as Theorem~\ref{thm:p02-w02c-p02-025-reduction} states.

If \(\delta_P(4,1)\) is changed to \((1,0,R)\), the computation halts in state
\(1\) at time \(3\), and \(\rho^{\mathrm B}_P\) becomes unsatisfiable: by
Lemma~\ref{thm:p02-w02c-p02-025-simulation}, every model contains an element
\(\tau_3\) with \(I_1\tau_3\), against \(\mathrm{NONSTOP}\).  If it is changed
to \((4,1,N)\), the program stays in state \(4\) forever; the intended model
on \(\mathbb N\) satisfies \(\rho^{\mathrm B}_P\), and every model contains
the pairwise distinct elements \(\tau_0,\tau_1,\ldots\), so it is infinite.
\end{hbexample}

\subsection{Limits of the construction}
\label{sec:p02-w02c-p02-025-limits}

\begin{hbwarning}{def:p02-w02c-p02-025-limits}{Where the construction needs its hypotheses}
\begin{enumerate}
\item \emph{Admissibility is used.}  Without (N2) the two regions collide.
  Let \(\delta_P(0,a)=(2,a,R)\) for both letters: the program halts in state
  \(2\) at time \(1\) with the head on cell \(1\), but it moves right at time
  \(0\) with \(h_0=0>0-2\).  In a model whose element \(\tau_0\) has the witness
  \(\tau_1\), the \(d=R\) clause at \(x=\tau_0\), \(u=\tau_1\), \(y=\tau_0\)
  records the move as \(H(\tau_1,\tau_0)\), which is a lower pair and so reads
  ``at time \(1\) the head is on cell \(0\)''; and the intended position,
  \(H(\tau_1,\tau_1)\), is excluded by \(\mathrm{OFFDIAG}\).  The proof of
  Lemma~\ref{thm:p02-w02c-p02-025-simulation} uses (N2) at exactly this point,
  and Lemma~\ref{thm:p02-w02c-p02-025-intended} uses it for \(\mathrm{OFFDIAG}\).
  The slowing step of Definition~\ref{def:p02-w02c-p02-025-simulation} is what
  makes every simulating program admissible.
\item \emph{The monadic vocabulary is not bounded.}  The formula
  \(\rho^{\mathrm B}_e\) has one monadic symbol for each state of \(P_e\), and
  the number of states grows with \(M_e\).  The binary vocabulary is fixed at
  three symbols.  Reducing the monadic symbols as well is the business of the
  source's Chapter~3, which starts from this class
  \cite[Sect.~3.2, pp.~115--116]{BorgerGradelGurevich1997}.
\item \emph{The finite models are not small in the size of the machine.}  The
  finite model for an accepting machine has one element per step of
  \(P_e\), and one step of \(M_e\) costs \(P_e\) between \(2(2k-1)\) and
  \(2(3k-2)\) steps, \(k\) the block length.
\item \emph{Negation is used.}  \(\beta_P\) uses \(\neg T_0\) for the letter
  \(1\) and negated \(K\)- and \(H\)-atoms in its premises.  The source's later
  refinements to Krom and Horn matrices
  \cite[Sect.~2.1.2, Exercise~2.1.37, p.~37]{BorgerGradelGurevich1997} are not
  obtained here.
\end{enumerate}
\end{hbwarning}

\subsection{Relationships and sources}
\label{sec:p02-w02c-p02-025-relationships}

The formula \(\rho_{M,\prec}\) of Definition~\ref{def:p02-w02-p02-036-refined}
and the formula \(\rho^{\mathrm B}_P\) give the same reduction property
(compare Theorem~\ref{thm:p02-w02-p02-036-refined} with
Theorem~\ref{thm:p02-w02c-p02-025-reduction}); they differ in what they spend.
\(\rho_{M,\prec}\) works for any Turing machine directly, with a binary symbol
for each tape letter, an explicit successor relation \(S\) and the prefix
\(\forall x\,\exists u\,\exists v\,\forall y\) in its infinity axiom.
\(\rho^{\mathrm B}_P\) uses one tape symbol, no successor relation, a single
\(\forall\exists\forall\) block, and needs the detour through binary admissible
programs.  Either family yields Trakhtenbrot's theorem in the inseparable form
of Theorem~\ref{thm:p02-w02-p02-036-inseparable}, since
\(\mathcal B\subseteq\mathrm{FO}\).  The register-machine version of the same
refinement, with Skolem normal form \(\forall x\,\forall y\), is the source's
Exercise~2.1.32, treated in \cref{sec:p02-w02b-p02-031-register}.

Sources: the exercise, its hint and the remark
\cite[Sect.~2.1.2, Exercise~2.1.31 and Remark, pp.~35--36]{BorgerGradelGurevich1997},
which credits the exercise to B\"uchi's 1962 paper on Turing machines and the
Entscheidungsproblem (the source's reference~[64]); the formulas
\(\mathrm{START}\), \(\mathrm{NONSTOP}\) and (2.1)--(2.3)
\cite[Sect.~2.1.1, pp.~20--22]{BorgerGradelGurevich1997}; the halting problems
\(H_1\), \(H_2\), \(H_\infty\) and the reduction property
\cite[Sect.~2.1.2, p.~34]{BorgerGradelGurevich1997}; the relativized infinity
axiom \(\prec_{K,M}\) \cite[Sect.~2.1.2, p.~35]{BorgerGradelGurevich1997};
B\"uchi formulas and the class \([\exists\land\forall\exists\forall,(\omega,3)]\)
\cite[Sect.~3.2, pp.~115--116]{BorgerGradelGurevich1997}.  The block
simulation, the slowing step, the three-region reading of \(H\) and \(T_0\), and
the conjuncts \(\mathrm{BLANK}\), \(\mathrm{OFFDIAG}\), \(\mathrm{MOVE}\),
\(\mathrm{ONEHEAD}\) and \(\mathrm{Act}\,x\to Kxu\) are this handbook's
completion of the exercise; the source leaves them to the reader.
